\documentclass[10pt,a4paper]{amsart}
\usepackage[margin=19mm]{geometry}
\usepackage{amsmath,amssymb,mathtools}
\usepackage[scr=rsfs,cal=euler]{mathalfa}
\usepackage{xfrac}
\usepackage[unicode,hidelinks,bookmarks=false]{hyperref}
\newcommand{\ie}{i.e.\ }
\newcommand{\cf}{cf.\ }
\newcommand{\resp}{respectively\ }
\newcommand{\Subsection}{\subsection}
\providecommand{\nobreakdash}{\nobreak\hbox{-}}
\setlength{\emergencystretch}{2em}
\multlinegap0pt
\usepackage{bm}
\usepackage{bbm}
\usepackage{dsfont}
\usepackage{xspace}
\usepackage{cancel}
\usepackage{mathtools}
\usepackage{cleveref}
\usepackage{amsthm}
\makeatletter
\@ifundefined{lemma}{\newtheorem{lemma}{Lemma}}{}
\makeatother
\title[On the Superconvergence of ESFR Schemes]{On the Superconvergence of ESFR Schemes}
\author{Mathias Dufresne-Pich\'e, Siva Nadarajah}
\date{Source: arXiv:2510.24111v1 (CC BY 4.0)}
\begin{document}
\maketitle

\noindent\emph{Source and licence.} On the Superconvergence of ESFR Schemes. Version arXiv:2510.24111v1 (CC BY 4.0), distributed under the Creative Commons Attribution 4.0 International licence, as recorded in the arXiv licence metadata for this exact version and shown on the abstract page https://arxiv.org/abs/2510.24111v1. Full rights notice: licenses/arxiv-2510.24111-CC-BY-4.0.txt.\par\smallskip

\noindent\textit{Editorial scope.} Counted unit: the Lemma whose source label is \texttt{lem:rational\_approx} in the article (source0.tex lines 309--323, source label \texttt{lem:rational\_approx}), delivered together with the article's own complete original proof of it (source0.tex lines 324--378, one \texttt{proof} environment of 55 lines). No mathematics is written, repaired or re-derived: statement and proof are the article's own text line for line. Required context retained uncounted: eq:tildeH (lines 271--274). Every definition, display and label of those lines is retained unchanged. Omitted from this package and not redistributed: 1--270, 275--308, 379--719. Nothing inside a retained excerpt is omitted or altered: every retained body is delivered exactly as the article's lines are written, so no omission receipt of any kind accompanies this package. Citations: the retained excerpts carry no citation command at all, so no bibliography entry of the article is transcribed and no reference is redistributed. Reference adaptation: none was needed and none was made; every reference inside the delivered unit resolves inside it, so no unresolved reference or citation is produced. Printed slips kept literal: no printed word, symbol, hypothesis or number of the counted statement or of its proof was altered. Layout and class: the article's own class is \texttt{siamart250211} with packages including amssymb, amsmath, bm, amsfonts, graphicx, subfig, epstopdf, algorithmic, amsopn; the wrapper is the lane's pinned \texttt{amsart} editorial wrapper at 10pt on A4, which restates the theorem-like environments the retained text needs and adds the article's own macro definitions that the retained text needs (none), with extra wrapper packages (bm, bbm, dsfont, xspace, cancel, mathtools, cleveref). The delivered numbering is the unit's own. Assets: no retained span contains a figure environment, a graphics-inclusion command, a PDF-page inclusion, a table or a TikZ picture; no image file is copied into this package, the package declares no asset dependency and no image file is redistributed with it. Licence: version arXiv:2510.24111v1 of this article is distributed under the Creative Commons Attribution 4.0 International licence, as recorded in the arXiv licence metadata for this exact version and shown on the abstract page https://arxiv.org/abs/2510.24111v1. A full rights notice is delivered at licenses/arxiv-2510.24111-CC-BY-4.0.txt. Characters: every retained excerpt is pure ASCII, so the delivered engine, XeLaTeX with its default Latin Modern text font, reports no missing character.
\begin{equation}
    \tilde{\mathbf{H}}(\theta) :=(\mathbf{M} + \mathbf{K})\left(\mathbf{D} -\frac{1}{2}i\theta \mathbf{I}\right) +\mathbf{l}\left(\mathbf{l}-\mathbf{r}e^{-i\theta}\right)^T.
    \label{eq:tildeH}
\end{equation}

\begin{lemma}
\label{lem:rational_approx}
Let $\psi^{(k)}_j(\xi)$ denote the $k$th derivative of the $j$th Legendre polynomial evaluated at $\xi$ and let $q_j$ denote the $j$th diagonal entry of $(\mathbf{M}+\mathbf{K})^{-1}$ when expressed in the Legendre basis. The determinant of the modified von Neumann matrix $\tilde{\mathbf{H}}$ is given by
\begin{equation}
     \text{\textup{det}}(\tilde{\mathbf{H}}(\theta)) =
    (-1)^{p+1}\text{\textup{det}}(\mathbf{M}+\mathbf{K})
    \left(P(i\theta) - Q(i\theta)e^{-i\theta}\right),
    \label{eq:simp1}
\end{equation}
where $P \in \mathcal{P}^{p+1}$ and $Q \in \mathcal{P}^{p}$ are defined as
\begin{align}
    P(z) &:=2^{-(p+1)}z^{p+1}-\sum_{j,k=0}^p (-1)^j2^{k-p} q_j\psi^{(k)}_j(-1)z^{p-k}, \label{eq:approximant_P} \\
    Q(z) &:=-\sum_{j,k=0}^p (-1)^j2^{k-p} q_j\psi^{(k)}_j(1)z^{p-k}. \label{eq:approximant_Q}
\end{align}
\end{lemma}

\begin{proof}
    For convenience and without loss of generality, we will assume that $\tilde{\mathbf{H}}$ is constructed in the Legendre basis. Leveraging the dyadic product structure of $\tilde{\mathbf{H}}$, we start by applying the well-known matrix determinant lemma to \cref{eq:tildeH}, which yields
\begin{equation}
    \text{det}(\tilde{\mathbf{H}}(\theta)) =
    \text{det}(\mathbf{A}(\theta))
    \left(1 + \mathbf{l}^T\mathbf{A}(\theta)^{-1} \mathbf{l} - \mathbf{r}^T\mathbf{A}(\theta)^{-1}\mathbf{l}e^{-i\theta}\right),
    \label{eq:det_lemma}
\end{equation}
where
\begin{equation}
    \mathbf{A}(\theta) := (\mathbf{M} + \mathbf{K})\left(\mathbf{D} -\frac{1}{2}i\theta \mathbf{I}\right).
\end{equation}
Since $\mathbf{D}$ is strictly upper triangular, we can deduce
\begin{equation}
    \text{det}(\mathbf{A}(\theta))
    = 
    \text{det}(\mathbf{M} +\mathbf{K})\left(-\frac{1}{2}i\theta\right)^{p+1}.
\end{equation}
To determine $\mathbf{A}(\theta)^{-1}$, we expand it in terms of its von Neumann series. In particular, since $\mathbf{D}$ is nilpotent with $\mathbf{D}^{p+1} = 0$, the von Neumann series of this operator can be written as the finite sum
\begin{align}
    \mathbf{A}(\theta)^{-1}
    &= -\left(\frac{1}{2}i\theta \mathbf{I} - \mathbf{D}\right)^{-1}(\mathbf{M}+\mathbf{K})^{-1} \\
    &=
    -\sum_{k=0}^{p} 2^{k+1}(i\theta)^{-(k+1)} \mathbf{D}^k (\mathbf{M}+\mathbf{K})^{-1}.
    \label{eq:vonneumann_series}
\end{align}
Combining \cref{eq:det_lemma} and \cref{eq:vonneumann_series} yields
\begin{equation}
     \text{det}(\tilde{\mathbf{H}}(\theta)) =
    (-1)^{p+1}\text{det}(\mathbf{M}+\mathbf{K})
    \left(2^{-(p+1)}(i\theta)^{p+1} + \mathbf{l}^T\mathbf{B}(\theta) \mathbf{l} - \mathbf{r}^T\mathbf{B}(\theta)\mathbf{l}e^{-i\theta}\right),
\end{equation}
where
\begin{equation}
    \mathbf{B}(\theta) := \sum_{k=0}^{p} 2^{k-p}(i\theta)^{p-k} \mathbf{D}^k (\mathbf{M}+\mathbf{K})^{-1}.
\end{equation}
Letting $q_j$ denote the $j$th diagonal entry of $(\mathbf{M}+\mathbf{K})^{-1}$ we leverage the structure of $\mathbf{B}$ and the fact $\mathbf{M} + \mathbf{K}$ is diagonal to find that
\begin{align}
    \mathbf{l}^T\mathbf{B}(\theta)\mathbf{l}
    &=
    -\sum_{k=0}^{p} 2^{k-p}(i\theta)^{p-k} \mathbf{l}^T\mathbf{D}^k (\mathbf{M}+\mathbf{K})^{-1} \mathbf{l} \\
    &=
    -\sum_{j=0}^p\sum_{k=0}^{p} (-1)^j2^{k-p} q_j\psi^{(k)}_j(-1)(i\theta)^{p-k}.
    \label{eq:left}
\end{align}
Similarly, it can be seen that
\begin{align}
    \mathbf{r}^T\mathbf{B}(\theta)\mathbf{l}
    &=
    -\sum_{k=0}^{p} 2^{k-p}(i\theta)^{p-k} \mathbf{r}^T\mathbf{D}^k (\mathbf{M}+\mathbf{K})^{-1} \mathbf{l} \\
    &=
    -\sum_{j=0}^p\sum_{k=0}^{p} (-1)^j2^{k-p} q_j\psi^{(k)}_j(1)(i\theta)^{p-k}.
    \label{eq:right}
\end{align}
\end{proof}

\end{document}
